%%%%%%%%%%%%%%%%%%%%%%%%%%%%%%%%%
%
%       EXERCÍCIO
%
%%%%%%%%%%%%%%%%%%%%%%%%%%%%%%%%%

\ifdefstring{\atividade}{prova}{%
    \renewcommand{\valorquestao}{\ValorQ\ pontos}
}%

\begin{exercicioBanco}[\valorquestao]
Três candidatos disputam as eleições para o Governo do Estado. O candidato do partido de direita tem \(30\%\) da preferência eleitoral, o de centro tem \(30\%\) e o da esquerda, \(40\%\). Em sendo eleito, a probabilidade de dar, efetivamente, prioridade para a Educação e Saúde é de \(0.4, 0.6\) e \(0.9\), para os candidatos de direita, centro e esquerda, respectivamente.
%
\begin{enumerate}[(a)]
    \item Qual é a probabilidade de não ser dada prioridade  a essas áreas no próximo governo?
    %
    \item Se a área teve prioridade, qual a probabilidade do candidato de direita ter ganho a eleição?
\end{enumerate}
%
\end{exercicioBanco}

